\chapter{Antipattern nelle architetture a microservizi}

\section{Antipatterns e code smells}
Come già discusso nelle sezioni precedenti, i microservizi stanno guadagnando sempre più popolarità e importanza grazie alle loro caratteristiche. Per facilitare tale sviluppo, come riportato nel paragrafo~\ref{sec:par2_4}, sono state identificate delle soluzioni a problemi di progettazione ricorrenti, tali soluzioni prendono il nome di \textit{design patterns} \cite{CernyCatalogo}.

Analogamente, gli \textbf{anti-patterns} sono identificabili come pratiche o scelte di progettazione a problemi ricorrenti che, pur sembrando efficienti in una prima fase, producono successivamente conseguenze negative e compromettono di conseguenza la qualità complessiva del sistema \cite{CernyCatalogo}.

Potremmo quindi definirli, come soluzioni \textit{"non ottimali"} a problemi ricorrenti. Un bad smell (o code smell), invece, rappresenta un sospetto che qualcosa non sia corretto e si configura a tutti gli effetti come un segnale di allarme che fa trasparire eventuali problemi nel design, ma non fornisce soluzione specifica \cite{CernyCatalogo}.
Citando Fowler, un code smell è "come un segnale superficiale che spesso indica un problema più profondo nel sistema" \cite{fowler_codesmell}.
Potremmo quindi identificare gli anti-pattern come la causa profonda del problema mentre i bad smells indicano i sintomi osservabili. 

L'origine degli anti-pattern è dovuta a varie cause che non sono solo vincolate ad aspetti tecnici ma anche a situazioni culturali e puramente organizzative:
\begin{itemize}
    \item Pressioni dovute a scadenze del progetto che portano a decisioni frettolose per concludere entro i termini previsti. 
    \item Mancanza di esperienza e competenza da parte del team.
    \item Mancanza di comunicazione tra i membri del team o tra stakeholder e sviluppatori. 
    \item Requisiti poco chiari, incompleti o in continua evoluzione. 
    \item Eccessiva fiducia in soluzioni adottate precedentemente con conseguente riutilizzo, ma come detto in precedenza le soluzioni vanno adottate in base al contesto di sviluppo che varia a seconda dei casi. 
    \item Scarsa documentazione del sistema. 
    \item Mancanza di test o testing inadeguato. 
\end{itemize}

Gli anti-pattern hanno ovviamente un impatto sul sistema. Le 

conseguenze che possiamo riscontrare sono le seguenti:
\begin{itemize}
    \item Aumento dell'accoppiamento tra i servizi, 
    
    compromettendo l'indipendenza.
    \item Aumento della complessità complessiva del sistema. 
    \item Difficoltà nel mantenere e nel testare il sistema, dovuto alla presenza di dipendenze implicite e interazioni non gestite opportunamente. 
    \item Degrado delle prestazioni, spesso causato da comunicazioni eccessive tra microservizi o da scelte architetturali inefficienti.
    \item Perdita di scalabilità a causa dei servizi che diventano dei veri e propri colli di bottiglia  difficili da isolare e ottimizzare. 
\end{itemize}

Ad esempio, nel caso di \textit{cyclic dependency}, due microservizi finiscono per dipendere reciprocamente e il fallimento di uno si propaga inevitabilmente all’altro. Analogamente, con il \textit{shared database} più microservizi accedono allo stesso database, aumentando l’accoppiamento e riducendo l’autonomia dei servizi.

In questo elaborato si fa riferimento alla convenzione terminologica adottata da Tomas Cerny e Davide Taibi, nella quale il termine anti-pattern è impiegato come categoria ombrello comprendente anche i bad smells, per ragioni di semplificazione concettuale \cite{CernyCatalogo}.

\section{Catalogo degli anti-pattern}

Lo studio condotto da Cerny e Taibi~\cite{CernyCatalogo} ha permesso di definire un catalogo composto da 58 anti-pattern, suddivisi in cinque macro-categorie:
\begin{itemize}
  \item \textbf{Intra-service design};
  \item \textbf{Inter-service decomposition};
  \item \textbf{Service interaction};
  \item \textbf{Security};
  \item \textbf{Team organization}.
\end{itemize}

Sono analizzate nel dettaglio nelle sezioni successive le singole categorie.

\subsection{Intra-service design}

Gli anti-pattern appartenenti a questa categoria fanno riferimento a difetti di progettazione interna di un singolo servizio. In particolare, riguardano aspetti come la granularità del servizio, la qualità dell’interfaccia esposta e la coesione funzionale delle responsabilità implementate.

La categoria \textit{Intra-service design} è suddivisa, in generale, in tre sottocategorie: \textit{granularity}, \textit{service interface} e \textit{functional cohesion}~\cite{CernyCatalogo}.

\begin{itemize}
    \item \textit{Granularity}: riguarda la dimensione del servizio, includendo casi come \textit{Nano-service} e \textit{Mega-service}.
    \item \textit{Service interface}: riguarda problemi relativi all’interfaccia del microservizio, come \textit{CRUDY service}, \textit{Nobody home}, \textit{Data service} e \textit{No API-versioning}.
    \item \textit{Functional cohesion}: riguarda la chiarezza, la comprensibilità e la coesione funzionale del servizio, includendo casi come \textit{Low cohesive operation} e \textit{Ambiguous service}.
\end{itemize}

Gli anti-pattern appartenenti a questa categoria sono riportati in Tabella~\ref{tab:intra_service_design}.

\Needspace{18\baselineskip}

\begin{table}[H]
\centering
\caption{Anti-pattern della categoria \textit{Intra-service design}.}
\label{tab:intra_service_design}
\small
\renewcommand{\arraystretch}{1.15}
\setlength{\tabcolsep}{7pt}

\begin{tabularx}{\textwidth}{
>{\centering\arraybackslash}p{0.8cm}
>{\RaggedRight\arraybackslash}p{3.8cm}
>{\RaggedRight\arraybackslash}X
}
\toprule
\textbf{Id} & \textbf{Anti-pattern} & \textbf{Descrizione} \\
\midrule

\multicolumn{3}{l}{\textit{Granularity}} \\
\cmidrule(lr){1-3}

1 & Nano-service
& Servizio eccessivamente piccolo, il cui costo di comunicazione e manutenzione supera i benefici della separazione. \\

\addlinespace[0.3em]

2 & Mega-service
& Servizio di dimensioni eccessive, con troppe responsabilità e ridotta manutenibilità. \\

\addlinespace[0.6em]

\multicolumn{3}{l}{\textit{Service interface}} \\
\cmidrule(lr){1-3}

3 & CRUDY service
& Servizio che espone principalmente operazioni CRUD, senza una reale logica di business. \\

\addlinespace[0.3em]

4 & Nobody home
& Servizio deployato e funzionante, ma non utilizzato o invocato da altri componenti. \\

\addlinespace[0.3em]

5 & Data service
& Servizio semplice che espone principalmente metodi per l’accesso ai dati. \\

\addlinespace[0.3em]

6 & No API-versioning
& Mancanza di versionamento delle API, con rischio di breaking changes per i client. \\

\addlinespace[0.6em]

\multicolumn{3}{l}{\textit{Functional cohesion}} \\
\cmidrule(lr){1-3}

7 & Whatever types
& Uso di tipi di dati generici che rendono ambiguo il contratto dell’interfaccia. \\

\addlinespace[0.3em]

8 & Low cohesive operation
& Operazioni con bassa coesione raggruppate nello stesso servizio. \\

\addlinespace[0.3em]

9 & Ambiguous service
& Servizio con responsabilità non ben definite o funzionalità sovrapposte. \\

\bottomrule
\end{tabularx}
\end{table}
\normalsize

\Needspace{20\baselineskip}
\subsection{Inter-service decomposition}

Gli anti-pattern appartenenti a questa categoria riguardano errori nella decomposizione del sistema in servizi e nelle relazioni strutturali tra essi. Nel dettaglio si basano su scelte errate nella suddivisione del sistema in microservizi, accoppiamenti stretti o violazione dei principi di modularità che caratterizzano architetture di questo genere.

La categoria \textit{Inter-service decomposition} include quindi problematiche che emergono non tanto all’interno del singolo servizio, quanto nel modo in cui i servizi vengono separati, organizzati e integrati tra loro~\cite{CernyCatalogo}.


Gli anti-pattern appartenenti a questa categoria sono riportati in Tabella~\ref{tab:inter_service_decomposition}.


\begin{table}[H]
\centering
\caption{Anti-pattern della categoria \textit{Inter-service decomposition}.}
\label{tab:inter_service_decomposition}
\small
\renewcommand{\arraystretch}{1.15}
\setlength{\tabcolsep}{7pt}

\begin{tabularx}{\textwidth}{
>{\centering\arraybackslash}p{0.8cm}
>{\RaggedRight\arraybackslash}p{3.8cm}
>{\RaggedRight\arraybackslash}X
}
\toprule
\textbf{Id} & \textbf{Anti-pattern} & \textbf{Descrizione} \\
\midrule

10 & Transactional integration
& Transazioni che si estendono oltre i confini del singolo servizio, riducendone l’autonomia. \\

\addlinespace[0.3em]

11 & Co-change coupling
& Servizi che devono essere modificati frequentemente insieme a causa di un forte accoppiamento. \\

\addlinespace[0.3em]

12 & Duplicated services
& Servizi multipli che implementano funzionalità simili o sovrapposte. \\

\addlinespace[0.3em]

13 & Microservice greedy
& Creazione eccessiva di microservizi senza chiari vantaggi misurabili. \\

\addlinespace[0.3em]

14 & Service chain
& Lunga catena di chiamate tra servizi, con conseguente aumento di latenza e complessità. \\

\addlinespace[0.3em]

15 & Hub-like dependency
& Servizio che diventa un punto centrale di dipendenza per molti altri servizi. \\

\addlinespace[0.3em]

16 & Cyclic dependency
& Dipendenze circolari tra servizi, che creano problemi di deployment e manutenzione. \\

\addlinespace[0.3em]

17 & Chatty service
& Comunicazione eccessiva e troppo fine-grained tra servizi, con effetti negativi sulle prestazioni. \\

\addlinespace[0.3em]

18 & Shared persistency
& Più servizi accedono allo stesso database, generando accoppiamento a livello dei dati. \\

\addlinespace[0.3em]

19 & Sand pile
& Sistema composto da molti nano-service che condividono dati comuni. \\

\addlinespace[0.3em]

20 & Shared libraries
& Condivisione di librerie runtime tra microservizi, con possibile compromissione dell’indipendenza. \\

\addlinespace[0.3em]

21 & Wrong cuts
& Decomposizione basata su layer tecnici invece che sul dominio di business. \\

\addlinespace[0.3em]

22 & Knot service
& Servizi fortemente accoppiati tramite integrazioni point-to-point. \\

\addlinespace[0.3em]

23 & Scattered parasitic functionality
& Funzionalità distribuite su più servizi, con responsabilità sovrapposte. \\

\bottomrule
\end{tabularx}
\end{table}
\normalsize

\subsection{Service interaction}

Gli anti-pattern appartenenti a questa categoria riguardano problemi nei meccanismi di comunicazione tra i servizi. Tale categoria di problemi può compromettere fattori importanti per il sistema come: la scalabilità, la manutenibilità e la capacità di evoluzione.

La categoria \textit{Service interaction} include quindi difetti legati al modo in cui i servizi si invocano, scambiano messaggi, gestiscono endpoint, fallimenti e informazioni operative~\cite{CernyCatalogo}.

Gli anti-pattern appartenenti a questa categoria sono riportati in Tabella~\ref{tab:service_interaction}.

\begin{table}[H]
\centering
\caption{Anti-pattern della categoria \textit{Service interaction}.}
\label{tab:service_interaction}
\small
\renewcommand{\arraystretch}{1.15}
\setlength{\tabcolsep}{7pt}

\begin{tabularx}{\textwidth}{
>{\centering\arraybackslash}p{0.8cm}
>{\RaggedRight\arraybackslash}p{3.8cm}
>{\RaggedRight\arraybackslash}X
}
\toprule
\textbf{Id} & \textbf{Anti-pattern} & \textbf{Descrizione} \\
\midrule

24 & ESB misuse
& Uso improprio di un Enterprise Service Bus come punto centrale di integrazione, con rischio di bottleneck. \\

\addlinespace[0.3em]

25 & On-line only
& Sistema che evita completamente l’elaborazione batch, anche quando questa sarebbe appropriata. \\

\addlinespace[0.3em]

26 & Empty messages
& Scambio di messaggi vuoti, utilizzati come semplici segnali tra componenti. \\

\addlinespace[0.3em]

27 & Use of business logic in communication
& Logica di business incorporata nello strato di comunicazione tra servizi. \\

\addlinespace[0.3em]

28 & Hardcoded endpoint
& Indirizzi ed endpoint hardcodati nel codice invece dell’utilizzo di meccanismi come il service discovery. \\

\addlinespace[0.3em]

29 & No API-gateway
& Mancanza di un gateway unificato, che costringe i client a chiamare direttamente i singoli servizi. \\

\addlinespace[0.3em]

30 & Wobbly service interactions
& Interazioni tra servizi che non gestiscono adeguatamente i fallimenti, riducendo la resilienza del sistema. \\

\addlinespace[0.3em]

31 & Timeout
& Timeout fissi che non considerano adeguatamente i diversi contesti operativi. \\

\addlinespace[0.3em]

32 & No health check
& Mancanza di endpoint o meccanismi per verificare lo stato di salute dei servizi. \\

\bottomrule
\end{tabularx}
\end{table}
\normalsize

\Needspace{20\baselineskip}
\subsection{Security}

Gli anti-pattern di questa categoria riguardano principalmente problematiche inerenti alla sicurezza dei microservizi.

La categoria \textit{Security} riguarda, quindi, la protezione dei sistemi attraverso tre meccanismi fondamentali: autenticazione, autorizzazione e cifratura. L’autenticazione è il processo di verifica dell’identità di un utente; l’autorizzazione determina quali risorse o operazioni sono accessibili a un utente autenticato; la cifratura consente invece di trasformare un testo leggibile in un formato illeggibile, al fine di proteggerne il contenuto~\cite{CernyCatalogo}.

Gli anti-pattern appartenenti a questa categoria sono riportati in Tabella~\ref{tab:security}.

\begin{table}[H]
\centering
\caption{Anti-pattern della categoria \textit{Security}.}
\label{tab:security}
\small
\renewcommand{\arraystretch}{1.15}
\setlength{\tabcolsep}{7pt}

\begin{tabularx}{\textwidth}{
>{\centering\arraybackslash}p{0.8cm}
>{\RaggedRight\arraybackslash}p{3.8cm}
>{\RaggedRight\arraybackslash}X
}
\toprule
\textbf{Id} & \textbf{Anti-pattern} & \textbf{Descrizione} \\
\midrule

33 & Unauthenticated traffic
& Traffico API non autenticato tra servizi o da sistemi esterni. \\

\addlinespace[0.3em]

34 & Multiple user authentication
& Presenza di molteplici punti di accesso per l’autenticazione utente, con aumento della superficie d’attacco. \\

\addlinespace[0.3em]

35 & Publicly accessible microservices
& Microservizi direttamente accessibili da client esterni senza adeguata protezione. \\

\addlinespace[0.3em]

36 & Unnecessary privileges
& Privilegi eccessivi concessi ai microservizi rispetto a quanto effettivamente necessario. \\

\addlinespace[0.3em]

37 & Insufficient access control
& Controlli di accesso mancanti o inadeguati sui microservizi. \\

\addlinespace[0.3em]

38 & Centralized authorization
& Autorizzazione centralizzata invece di una gestione più granulare a livello di singolo servizio. \\

\addlinespace[0.3em]

39 & Non-secured service communications
& Comunicazioni tra servizi non cifrate o non autenticate. \\

\addlinespace[0.3em]

40 & Non-encrypted data exposure
& Dati sensibili esposti in chiaro nello storage o durante il trasporto. \\

\addlinespace[0.3em]

41 & Own crypto code
& Implementazione personalizzata di algoritmi crittografici invece dell’utilizzo di librerie standard. \\

\addlinespace[0.3em]

42 & Hardcoded secrets
& Segreti, come password o API key, hardcodati nel codice o nelle configurazioni. \\

\bottomrule
\end{tabularx}
\end{table}
\normalsize

\Needspace{20\baselineskip}
\subsection{Team organization}

Gli anti-pattern appartenenti a questa categoria riguardano problemi organizzativi e operativi legati al modo in cui i team progettano, sviluppano, rilasciano e gestiscono i microservizi. A differenza delle categorie precedenti, il focus non è solo sulla struttura tecnica del software, ma anche sulle pratiche adottate durante il ciclo di vita del sistema.

La categoria \textit{Team organization} include quindi difetti legati sia alle scelte dei team di sviluppo sia alla gestione operativa del sistema, comprese eventuali violazioni delle pratiche DevOps~\cite{CernyCatalogo}.

Gli anti-pattern appartenenti a questa categoria sono riportati in Tabella~\ref{tab:team_organization}.

\begin{table}[H]
\centering
\caption{Anti-pattern della categoria \textit{Team organization}.}
\label{tab:team_organization}
\small
\renewcommand{\arraystretch}{1.15}
\setlength{\tabcolsep}{7pt}

\begin{tabularx}{\textwidth}{
>{\centering\arraybackslash}p{0.8cm}
>{\RaggedRight\arraybackslash}p{3.8cm}
>{\RaggedRight\arraybackslash}X
}
\toprule
\textbf{Id} & \textbf{Anti-pattern} & \textbf{Descrizione} \\
\midrule

\multicolumn{3}{l}{\textit{Development}} \\
\cmidrule(lr){1-3}

43 & Shiny nickel
& Adozione di nuove tecnologie solo perché di tendenza, senza reali benefici per il sistema. \\

\addlinespace[0.3em]

44 & Golden hammer
& Uso della stessa tecnologia familiare per ogni problema, anche quando non appropriata. \\

\addlinespace[0.3em]

45 & Lack of communication standards
& Assenza di standard comuni per API e formati di messaggio tra team. \\

\addlinespace[0.3em]

46 & Too many standards
& Eccessiva eterogeneità di linguaggi, protocolli e framework nell’organizzazione. \\

\addlinespace[0.3em]

47 & Inadequate techniques support
& Uso di tecniche inadeguate, che non supportano efficacemente lo sviluppo di microservizi. \\

\addlinespace[0.3em]

48 & Single layer team
& Team organizzati per layer tecnici invece che per servizio o dominio. \\

\addlinespace[0.3em]

49 & No legacy
& Tendenza a eliminare completamente i sistemi legacy invece di integrarli progressivamente. \\

\addlinespace[0.3em]

50 & Data-driven migration
& Migrazione simultanea di dati e funzionalità durante la scomposizione del sistema. \\

\addlinespace[0.3em]

51 & Big bang
& Riscrittura completa del sistema in una sola volta invece che tramite un processo incrementale. \\

\addlinespace[0.3em]

52 & Multiple service instances per host
& Deployment di più istanze di servizi sullo stesso host, con conseguente condivisione delle risorse. \\

\addlinespace[0.6em]

\multicolumn{3}{l}{\textit{Operation}} \\
\cmidrule(lr){1-3}

53 & No CI/CD
& Mancanza di integrazione e deployment continui, con processi prevalentemente manuali. \\

\addlinespace[0.3em]

54 & Manual configuration
& Configurazione manuale invece di automazione e strumenti di configuration management. \\

\addlinespace[0.3em]

55 & Insufficient monitoring
& Monitoraggio inadeguato delle prestazioni e dei failure dei servizi. \\

\addlinespace[0.3em]

56 & Dismiss documentation
& Documentazione insufficiente delle API e dell’architettura. \\

\addlinespace[0.3em]

57 & Insufficient message traceability
& Messaggi privi di metadati sufficienti per tracciarne origine e percorso. \\

\addlinespace[0.3em]

58 & Local logging
& Log scritti solo in storage locale invece che inviati a sistemi centralizzati. \\

\bottomrule
\end{tabularx}
\end{table}
\normalsize